\documentclass[11pt,aspectratio=169]{beamer}
\usetheme{Boadilla}
\usefonttheme{professionalfonts}
\setbeamertemplate{navigation symbols}{}

\usepackage[T1]{fontenc}
\usepackage{lmodern}
\usepackage{amsmath,amssymb,mathtools}
\usepackage{booktabs}
\usepackage{tabularx}
\usepackage{array}
\usepackage{graphicx}
\usepackage[percent]{overpic}
\usepackage{xcolor}
\usepackage[authoryear,round]{natbib}
\usepackage{bibentry}
\bibliographystyle{apalike}
\setbeamertemplate{bibliography item}{}
\makeatletter
\def\hyper@natlinkstart#1{}
\def\hyper@natlinkend{}
\def\hyper@natlinkbreak#1#2{#1}
\makeatother

\graphicspath{{figs/}}

% Section headline for the audience arc.
\setbeamertemplate{headline}{%
  \hbox{}\vskip2pt\hfill\textcolor{gray!70}{\scriptsize\insertsectionhead}\hspace{1em}\vskip-6pt%
}
\setbeamertemplate{section in toc}[sections numbered]
\setbeamertemplate{frametitle}[default][left]
\setbeamerfont{framesubtitle}{size=\small}

\newcommand{\figplaceholder}[1]{\begin{center}\fbox{\begin{minipage}[c][0.32\textheight][c]{0.82\textwidth}\centering\large #1\end{minipage}}\end{center}}
\newcommand{\shortfigplaceholder}[1]{\begin{center}\fbox{\begin{minipage}[c][0.22\textheight][c]{0.82\textwidth}\centering\large #1\end{minipage}}\end{center}}
\newcommand{\compactfigplaceholder}[1]{\begin{center}\fbox{\begin{minipage}[c][0.10\textheight][c]{0.82\textwidth}\centering\normalsize #1\end{minipage}}\end{center}}
\newcommand{\R}{\mathcal{R}}
\newcommand{\E}{\mathcal{E}}
\newcommand{\dd}{\,\mathrm{d}}
\newcommand{\good}[1]{\textcolor{blue!70!black}{#1}}
\newcommand{\bad}[1]{\textcolor{red!70!black}{#1}}
\newcommand{\SAR}{speed asymmetry ratio}
% Inline slide navigation (end of sentence / bullet).
\newcommand{\app}[1]{\allowbreak\hspace{0.25em}\hyperlink{#1}{\beamergotobutton{details}}}
\newcommand{\back}[1]{\allowbreak\hspace{0.25em}\hyperlink{#1}{\beamerreturnbutton{back}}}
% Appendix / frame footer: bottom-right, not centered.
\newcommand{\navback}[1]{\par\vspace{0.35em}\noindent\hfill\hyperlink{#1}{\beamerreturnbutton{back}}}
\newcommand{\navfwd}[1]{\par\vspace{0.35em}\noindent\hfill\hyperlink{#1}{\beamergotobutton{details}}}
\newcommand{\navbackfwd}[2]{\par\vspace{0.35em}\noindent\hfill\hyperlink{#1}{\beamerreturnbutton{back}}\;\hyperlink{#2}{\beamergotobutton{details}}}

\title[Directional Congestion]{Directional Congestion}
\subtitle{Measurement, Spatial Equilibrium, and Tidal-Lane Reallocation}
\author[Duan, Feng, Gu]{Yige Duan (SJTU) $\quad$ Xiaoruo Feng (PKU) $\quad$ Yizhen Gu (PKU)}
\date[SMUE Wuhan, 2026]{SMUE Wuhan, 2026}

\begin{document}

\begin{frame}
\titlepage
\end{frame}

% =====================================================================
% Opener: roadmap followed by examples of directional congestion.
% =====================================================================
\begin{frame}{Roadmap}
\footnotesize\tableofcontents
\end{frame}

\begin{frame}{Directional Congestion in Cities}

\centering
% All four photos center-cropped to 3:2 for a uniform grid.
\begin{minipage}[t]{0.32\textwidth}\centering
\includegraphics[width=0.8\linewidth]{cityBeijing_c.jpg}\\[0.02em]{\scriptsize Beijing}
\end{minipage}\hspace{0.02\textwidth}%
\begin{minipage}[t]{0.32\textwidth}\centering
\includegraphics[width=0.8\linewidth]{cityWuhan_c.jpg}\\[0.02em]{\scriptsize Wuhan}
\end{minipage}

\vspace{0.25em}
\begin{minipage}[t]{0.32\textwidth}\centering
\includegraphics[width=0.8\linewidth]{cityShenzhen2_c.jpg}\\[0.02em]{\scriptsize Shenzhen}
\end{minipage}\hspace{0.02\textwidth}%
\begin{minipage}[t]{0.32\textwidth}\centering
\includegraphics[width=0.8\linewidth]{cityKualaLumpur_c.jpg}\\[0.02em]{\scriptsize Kuala Lumpur}
\end{minipage}
\end{frame}

% =====================================================================
\section{Introduction}
% =====================================================================
\begin{frame}[label=motivation]{Motivation}
\begin{columns}[T,onlytextwidth]
\column{0.65\textwidth}
\begin{itemize}\small
  \item \textbf{Directional congestion is widely observed on urban roads, particularly during peak hours.}
  \begin{itemize}
    \item Commuting traffic is shaped by where workers live and where jobs are located.
  \end{itemize}
  \item \textbf{The literature often ignores this asymmetry in congestion.}
  \begin{itemize}
    \item Examples: QSM applications usually enter travel frictions through bilateral or route-level costs \citep{ahlfeldt2015,allenarkolakis2022}.
    \item Except: transport studies discuss peak-direction traffic \citep{vickrey1969,wolshon2006}.
  \end{itemize}
  \item \textbf{Tidal lanes are proposed as a low-cost way to address it.}
  \begin{itemize}
    \item However, little is known about the short-run and long-run effects of tidal lanes.
  \end{itemize}
\end{itemize}
\column{0.34\textwidth}
\centering
\href{run:./A6_XiaoruoFeng_DirectionalCongestion_video.mov}{%
  \includegraphics[width=\linewidth,height=0.70\textheight,keepaspectratio]{TidalVedio/pic.jpg}}\\[3pt]
{\scriptsize \good{Click to play}\app{appvideofacts}}
\end{columns}
\end{frame}

\begin{frame}{This Paper}
\small
\begin{itemize}
  \item \textbf{Measure directionality, using speed data at both the trip and the road segment level.}
  \begin{itemize}\footnotesize
    \item National trip-level: flow-weighted median \SAR\ is \good{0.878}.
    \item Within-city road segments: \SAR\ $\leq1/3$ flags \good{179} Beijing AM two-way centerlines (\good{4.9\%}, about \good{213 km}).
  \end{itemize}
  \item \textbf{Quantify what reciprocal symmetry misses in a spatial GE model.}
  \begin{itemize}\footnotesize
    \item \good{$56.3\%$} of passages take the slower direction; symmetry understates travel time by \good{about $5\%$}.
    \item Symmetry understates peak travel time from outer residential areas and therefore central residential fundamentals.
  \end{itemize}
  \item \textbf{Evaluate counterfactuals under endogenous congestion GE.}
  \begin{itemize}\footnotesize
    \item Tidal top-100 shifts one lane toward the slow AM direction on each treated corridor.
    \item \textbf{ST:} total AM-peak time falls by \good{$1{,}031$ hours} and welfare rises by \good{$0.87\%$}; \textbf{LT sorting:} the time gain reverses, leaving total time \bad{$352$ hours above baseline}, while the welfare gain narrows to \good{$0.32\%$}.
    \item Capacity-expansion benchmark builds an additional \good{$1{,}100$ lane-km}, while tidal lanes add \good{zero}.
  \end{itemize}
\end{itemize}
\end{frame}
\note{\textbf{Talk track, 55 seconds.} The paper makes three contributions. First, it measures directional congestion at the national trip level and on Beijing road segments. Second, it quantifies what reciprocal symmetry misses. The more congested direction carries 56.3 percent of AM link passages, so symmetry understates experienced travel time by about 5 percent. It also understates peak travel time from outer residential areas and therefore central residential fundamentals. Third, the model evaluates tidal reallocation and capacity expansion under endogenous congestion and long-run sorting.}

\begin{frame}{Literature}
\begin{itemize}\small
  \item \textbf{Directional congestion}
  \begin{itemize}\footnotesize
    \item Peak-period bottlenecks, measured urban speeds, and reversible-lane operations \citep{vickrey1969,arnott1993,couture2018,akbar2023,wolshon2006}.
    \item \good{This paper measures directionality at national trip level and road-segment level using high-frequency speed data.}
  \end{itemize}
  \item \textbf{Spatial equilibrium}
  \begin{itemize}\footnotesize
    \item Economic geography and quantitative spatial models \citep{krugman1991,eaton2002,allen2014,ahlfeldt2015}.
    \item \good{This paper applies spatial GE with directional congestion and analyze tidal lane counterfactuals}
  \end{itemize}
  \item \textbf{Transportation policy and congestion}
  \begin{itemize}\footnotesize
    \item Demand-side policies: congestion pricing and equilibrium sorting \citep{vickrey1969,kreindler2024,barwick2024,yang2020}.
    \item Supply-side policies: road capacity, fundamental law, and optimal networks \citep{duranton2011,fajgelbaum2020}.
    \item Public transit as a substitute: subways relieve road congestion \citep{gu2021}.
    \item \good{This paper evaluates tidal lanes as a directed capacity-reallocation policy.}
  \end{itemize}
\end{itemize}
\end{frame}

\begin{frame}[label=data]{Data}
\small
\begin{columns}[T,onlytextwidth]
\column{0.5\textwidth}
\begin{itemize}
  \item \textbf{Directed speed.}
  \begin{itemize}\footnotesize
    \item National trip-level routes across Chinese cities.
    \item Road-segment panel at 10-minute frequency, over one week in 2025.
  \end{itemize}
  \item \textbf{Commuting flows.}
  \begin{itemize}\footnotesize
    \item Baidu 2022 home--work commuting flow data.
  \end{itemize}
  \item \textbf{Road characteristics.}
  \begin{itemize}\footnotesize
    \item Road geometry, functional class, and lane records.
  \end{itemize}
\end{itemize}
\column{0.50\textwidth}
\centering
\includegraphics[width=\linewidth,height=0.76\textheight,keepaspectratio]{speed_eg.png}
\end{columns}
\navfwd{appmeasurementrouting}
\end{frame}

\begin{frame}[label=nattrip]{Directionality Appears at National and Road-Segment Scales}
\begin{itemize}\footnotesize
  \item \textbf{Speed asymmetry ratio.} The lower speed divided by the higher speed across a pair's two directions. It ranges from 0 to 1; 1 means equal speeds, and lower values indicate stronger directional imbalance.
\end{itemize}
\vspace{0.20em}
\begin{columns}[T,onlytextwidth]
\column{0.605\textwidth}
\centering\scriptsize \textbf{National trip-level directionality}\par
\vspace{0.35em}
\begin{minipage}[t]{0.49\linewidth}\centering
\includegraphics[width=\linewidth,height=0.30\textheight,keepaspectratio]{pat_national_ratio_hist.png}
\end{minipage}\hspace{0.01\linewidth}%
\begin{minipage}[t]{0.49\linewidth}\centering
\includegraphics[width=\linewidth,height=0.30\textheight,keepaspectratio]{pat_congestion_am_scatter.png}
\end{minipage}
\vspace{0.10em}
\begin{itemize}\scriptsize\setlength\itemsep{0.45em}
  \item Across the top 100 congested cities, the flow-weighted median is \good{0.878}; \good{58\%} of pair-flow occurs on pairs with ratios $\leq0.90$.
  \item More congested cities tend to have lower ratios and stronger asymmetry ($r=\bad{-0.48}$, $n=99$).
\end{itemize}
\column{0.385\textwidth}
\centering\scriptsize \textbf{Within-city road-segment directionality}\par
\vspace{0.35em}
\includegraphics[width=\linewidth,height=0.30\textheight,keepaspectratio]{pat_bj_centerline_map.png}
\vspace{0.10em}
\begin{itemize}\scriptsize\setlength\itemsep{0.45em}
  \item Median \SAR\ is \good{0.884} among 3{,}670 two-way centerlines.
  \item \SAR\ $\leq1/3$ flags \good{179} candidate corridors, about \good{213 km}.
\end{itemize}
\end{columns}
\end{frame}

\begin{frame}[label=jobhousing]{Directionality Arises from Jobs--Housing Mismatch}
\begin{columns}[T,onlytextwidth]
\column{0.48\textwidth}
\centering
\small\textbf{Across major cities}\par
\vspace{0.30em}
\begin{itemize}\footnotesize\setlength\itemsep{0.55em}
  \item \textbf{Widespread.} Every top-20 city has more AM commuting toward its employment core than away from it. The flow-weighted gap is \good{4.5 pp}.
\end{itemize}
\vspace{2.6em}
\begin{minipage}{0.92\linewidth}
\makebox[\linewidth][c]{\includegraphics[width=0.96\linewidth,height=0.38\textheight,keepaspectratio]{pat_jobs_housing_top20_gap_lollipop.pdf}}
\vspace{-0.35em}

{\tiny\raggedright\textit{Notes:} Core towns cover the first 30\% of workplace-side commuter mass; all other towns are peripheral.\par}
\end{minipage}

\column{0.50\textwidth}
\centering
\small\textbf{Beijing}\par
\vspace{0.30em}
\begin{itemize}\footnotesize\setlength\itemsep{0.55em}
  \item \textbf{Separation.} Jobs are more concentrated than residences: Duncan $D=\good{0.26}$ versus $0.004$ under random allocation. \app{appduncan}
  \item \textbf{Direction.} Peripheral-to-core commuting exceeds the reverse by \good{9.4 pp} in Beijing.
\end{itemize}
\vspace{0.35em}
\includegraphics[width=0.96\linewidth,height=0.34\textheight,keepaspectratio]{pat_jobs_housing_balance.png}
\end{columns}
\end{frame}

% =====================================================================
\section{Model}
% =====================================================================
% Keep the model in link-cost notation. Physical minutes m_kl enter only in
% the measurement and physical-time accounting slides, where t_kl=m_kl^(1/theta).

\begin{frame}{Model: Environment}
\begin{itemize}\small\setlength\itemsep{0.70em}
  \item Locations $i,j\in N$ host a fixed mass $\bar L$ of residents and jobs.
  \item A physical road is represented by two directed links, $(k,l)$ and $(l,k)$.
  \begin{itemize}\footnotesize\setlength\itemsep{0.30em}
    \item $t_{kl}^{ff}$ is the free-flow link cost on directed link $(k,l)$.
    \item $n_{kl}$ is the number of lanes assigned to that direction.
  \end{itemize}
  \item Realized link cost $t_{kl}$ is endogenous and directed.
  \begin{itemize}\footnotesize\setlength\itemsep{0.30em}
    \item $t_{kl}$ rises with traffic per lane, $\Xi_{kl}/n_{kl}$.
    \item $t_{kl}\neq t_{lk}$ when traffic differs across directions.
  \end{itemize}
  \item A tidal-lane policy shifts one lane across reciprocal directions while holding $n_{kl}+n_{lk}$ fixed.
\end{itemize}
\end{frame}
\note{\textbf{Talk track, 35 seconds.} A physical road becomes two directed links. Free-flow link cost and lanes are directed primitives. Realized link cost responds to traffic per lane, so the two directions can differ on the same road. A tidal lane changes directed capacity without expanding total road width.}

\begin{frame}{Individual Location Choice}
\begin{itemize}\small\setlength\itemsep{0.75em}
  \item An individual $\omega$ living in $i$, working in $j$, and taking route $r\in\R_{ij}$ obtains
  \vspace{-0.10em}\begin{footnotesize}\begin{equation}\label{eq:routecost}
    V_{ij,r}(\omega) = \left( \frac{u_i w_j}{C_{ij,r}} \right)\varepsilon_{ij,r}(\omega),
    \qquad
    C_{ij,r} = \prod_{\ell=1}^{K_r} t_{r_{\ell-1},r_\ell}
  \vspace{-0.35em}\end{equation}\end{footnotesize}
  \begin{itemize}\footnotesize\setlength\itemsep{0.30em}
    \item $u_i$ is amenity; $w_j$ is wage; $\varepsilon_{ij,r}$ is i.i.d. Fr\'echet.
    \item $C_{ij,r}$ is the route-specific travel cost for $r=[r_0,\ldots,r_{K_r}]$, with $r_0=i$ and $r_{K_r}=j$.
  \end{itemize}
  \vspace{0.55em}
  \item Residence, workplace, and route are chosen jointly:
  \vspace{-0.10em}\begin{equation}\label{eq:pi}
    \pi_{ij,r} = \frac{\left(u_i w_j/C_{ij,r}\right)^\theta}{\sum_{a}\sum_b\sum_{s\in\R_{ab}}\left(u_a w_b/C_{ab,s}\right)^\theta}.
  \end{equation}
\end{itemize}
\end{frame}
\note{\textbf{Talk track, 45 seconds.} The choice object is $(i,j,r)$. Route cost is the product of directed link costs along the route. Fr\'echet heterogeneity delivers a tractable probability for every residence, workplace, and route combination.}

\begin{frame}[label=bilateral]{Bilateral Commuting Costs and Flows}
\begin{itemize}\small
  \item Aggregating route costs gives the bilateral commuting cost:
  \vspace{-0.5em}\begin{equation}\label{eq:bilateral}
    \tau_{ij} = \left[\sum_{r\in\R_{ij}} C_{ij,r}^{-\theta}\right]^{-1/\theta}.
  \end{equation}
  \vspace{-0.5em}
  \item Summing route-specific flows gives residence-workplace commuting gravity:
  \vspace{-0.35em}\begin{equation}\label{eq:gravity}
    L_{ij}=\sum_{r\in\R_{ij}}L_{ij,r}
    =\tau_{ij}^{-\theta}\left(\frac{L_i^R}{\Pi_i}\right)\left(\frac{L_j^F}{P_j}\right),
    \quad L_i^R=\sum_jL_{ij},\quad L_j^F=\sum_iL_{ij}.
  \vspace{-0.8em}\end{equation}
  \begin{itemize}\footnotesize
    \item $\Pi_i$ and $P_j$ are origin and destination market-access terms.
    \item Directed link costs imply $\tau_{ij}\neq\tau_{ji}$ in general. \app{approute}
  \end{itemize}
\end{itemize}
\end{frame}
\note{\textbf{Talk track, 45 seconds.} The model first assigns probability to each route. Summing over routes produces one bilateral cost $\tau_{ij}$ and total OD flow $L_{ij}$. This gravity equation is route-inclusive because every directed link enters through $C_{ij,r}$.}

\begin{frame}{Amenity, Productivity, and Location Choices}
\begin{itemize}\small
  \item Local outcomes combine fundamentals with endogenous density:
  \vspace{-0.5em}\begin{equation}\label{eq:agglom}
    w_j = A_j = \bar A_j\left(L_j^F\right)^\alpha,
    \qquad
    u_i=\bar u_i\left(L_i^R\right)^\beta.
  \end{equation}
  \vspace{-0.4em}
  \item Residence ($R$) and workplace ($F$) shares satisfy
  \vspace{-0.2em}
  \begin{align}
  (l_i^R)^{1-\beta\theta} & = \chi \sum_j\tau_{ij}^{-\theta}\,\bar u_i^{\theta}\,\bar A_j^{\theta}\,(l_j^F)^{\alpha\theta} \label{eq:clear1} \\
  (l_j^F)^{1-\alpha\theta} & = \chi \sum_i\tau_{ij}^{-\theta}\,\bar u_i^{\theta}\,\bar A_j^{\theta}\,(l_i^R)^{\beta\theta}. \label{eq:clear2}
  \end{align}
  \vspace{-0.4em}
\end{itemize}
\end{frame}
\note{\textbf{Talk track, 45 seconds.} Fundamentals are $\bar u_i$ and $\bar A_j$. Employment density affects productivity through $\alpha$, and residential density affects amenities through $\beta$. Bilateral commuting costs connect these two location distributions.}

\begin{frame}{Feedback between Location Choice, Route Traffic, and Traversal Cost}
\begin{itemize}\footnotesize
  \item $\kappa^{kl}_{ij}$: \textbf{expected use of directed link $(k,l)$} by an $i\to j$ commuter.
  \vspace{-0.8em}\begin{equation}
    \kappa^{kl}_{ij}=\left(\frac{\tau_{ij}}{\tau_{ik}\,t_{kl}\,\tau_{lj}}\right)^\theta.
  \vspace{-0.6em}\end{equation}
  \item Directed link traffic aggregates OD flows across routes:
  \vspace{-0.5em}\begin{equation}\label{eq:Xi_c}
    \Xi_{kl}=\sum_i\sum_j\kappa^{kl}_{ij}L_{ij}.
  \vspace{-0.6em}\end{equation}
  \item Traffic per lane raises the directed link cost:
  \vspace{-0.5em}\begin{equation}\label{eq:c_Xi}
    t_{kl}=t_{kl}^{ff}\left(\frac{\Xi_{kl}}{n_{kl}}\right)^{\lambda/\theta}.
  \vspace{-0.6em}\end{equation}
  \begin{itemize}\scriptsize
    \item $\lambda$ is the physical time-flow elasticity; $\lambda/\theta$ follows from $t_{kl}=m_{kl}^{1/\theta}$.
  \end{itemize}
  \item Expected welfare is
  \vspace{-0.5em}\begin{equation}
    \bar W=\left[\sum_{i,j}\tau_{ij}^{-\theta}u_i^\theta w_j^\theta\right]^{1/\theta}.
  \vspace{-0.8em}\end{equation}
\end{itemize}
\end{frame}
\note{\textbf{Talk track, 60 seconds.} $\kappa_{ij}^{kl}$ maps one OD commuter into expected passages over a directed link. Summing over OD flows gives model-implied link traffic $\Xi_{kl}$. Traffic raises the directed link cost $t_{kl}$, which feeds back into location and route choices. This fixed point is the endogenous-congestion channel.}

\begin{frame}[label=equilibrium]{Equilibrium}
\begin{itemize}\footnotesize\setlength\itemsep{0.65em}
  \item Given fundamentals $\{\bar u_i,\bar A_i,t_{kl}^{ff},n_{kl}\}$ and parameters $(\theta,\alpha,\beta,\lambda)$, an equilibrium is
  \vspace{-0.10em}\begin{equation}
    \left\{L_i^R,L_j^F,L_{ij},t_{kl},\tau_{ij},\Xi_{kl},\Pi_i,P_j\right\}
  \vspace{-0.25em}\end{equation}
  such that
  \begin{itemize}\footnotesize\setlength\itemsep{0.26em}
    \item residence-workplace-route choices satisfy \eqref{eq:pi};
    \item bilateral flows satisfy gravity \eqref{eq:gravity} and location shares satisfy \eqref{eq:clear1}--\eqref{eq:clear2};
    \item amenities and productivity satisfy \eqref{eq:agglom};
    \item bilateral costs aggregate route-specific directed link costs through \eqref{eq:routecost} and \eqref{eq:bilateral};
    \item directed link traffic and traversal costs are mutually consistent \eqref{eq:Xi_c}--\eqref{eq:c_Xi}.
  \end{itemize}
\end{itemize}
\end{frame}
\note{\textbf{Talk track, 35 seconds.} Equilibrium closes the loop among location choice, route choice, link traffic, and directed link costs. The empirical section recovers the fundamentals and anchors the congestion equation to observed conditions before solving counterfactual equilibria.}

% =====================================================================
\section{Data and Estimation}\label{sec:quant}
% =====================================================================
\begin{frame}[label=estparams]{Estimation}
\small
\textbf{Data inputs}
\vspace{0.25em}

\scriptsize
\begin{tabularx}{0.92\textwidth}{@{}p{0.24\textwidth}X@{}}
\toprule
Object & Source \\
\midrule
$L_i^R,L_j^F,L_{ij}$ & Baidu 2022 commuting flow data \\
$t_{kl}^{\mathrm{obs}}$ & Baidu segment-level speed data: weekday AM peak average (07:00--09:00), May 19--25, 2025 \\
$t_{kl}^{ff}$ & Baidu segment-level speed data: night free-flow average (22:00--06:00, all days), May 19--25, 2025 \\
$\ell_{kl},n_{kl}$ & Road geometry and lane records \\
$(\mathcal N,\mathcal E)$ & Beijing 3 km grid and directed adjacent-grid network \\
\bottomrule
\end{tabularx}

\vspace{1.0em}
\textbf{Parameters}
\vspace{0.25em}

\scriptsize
\begin{tabularx}{0.92\textwidth}{@{}p{0.16\textwidth}p{0.15\textwidth}X@{}}
\toprule
Parameter & Value & Role / source \\
\midrule
$\theta$ & $6.83$ & cost elasticity (Ahlfeldt et al., 2015; Allen and Arkolakis, 2022) \\
$\alpha$ & $+0.05$ & production agglomeration \\
$\beta$ & $-0.10$ & residential crowding \\
$\lambda$ & $1.1$ & speed-flow elasticity (Barwick et al.\ 2024) \\
\bottomrule
\end{tabularx}
\navfwd{appparams}
\end{frame}
\note{\textbf{Talk track, 65 seconds.} The upper table maps the model objects to their data sources. Baidu 2022 supplies residence, workplace, and bilateral commuting flows. Observed directed link costs use weekday AM speeds from 07:00 to 09:00; free-flow link costs use the 22:00 to 06:00 average across all days. Road geometry and lanes determine link length and capacity, and the 3 km grid defines directed adjacency. The lower table reports the four calibrated parameters.}

\begin{frame}[label=linkmeasurement]{From Segment Speeds to Directed Link Costs}
\begin{center}
\begin{minipage}[t]{0.46\textwidth}
\centering
{\scriptsize One grid-link is a selected raw-road path}\par
\vspace{0.10em}
\includegraphics[width=0.98\linewidth,height=0.33\textheight,trim=0 130 0 18,clip,keepaspectratio]{figs/link_time_construction_actual.pdf}
\end{minipage}\hspace{0.035\textwidth}%
\begin{minipage}[t]{0.46\textwidth}
\centering
{\scriptsize Observed physical time on grid-links}\par
\vspace{0.10em}
\includegraphics[width=0.98\linewidth,height=0.33\textheight,keepaspectratio]{figs/app_edgeproj_adjacent_m_obs_slide.pdf}
\end{minipage}
\end{center}
\vspace{0.30em}
\begin{itemize}\footnotesize\setlength\itemsep{0.28em}
  \item \textbf{Pair-specific path.} For each adjacent grid pair $(k,l)$, select the directed raw-road path $p_{kl}$ that minimizes observed AM minutes; evaluate observed and free-flow speeds on that same path. \app{appmeasurementrouting}
  \item \textbf{Physical minutes.} Add raw-segment traversal minutes and off-road connector minutes $c_k^q,c_l^q$:
  \vspace{-0.40em}
  \begin{align}
    m_{kl}^{q}=c_{k}^{q}
    +60\sum_{a\in p_{kl}}\frac{\ell_a}{v_a^{q}}
    +c_{l}^{q},
    \qquad q\in\{obs,ff\}. \label{eq:linkminutes}
  \end{align}
  \vspace{-0.55em}
  \item \textbf{Model link cost.} $t_{kl}^{q}=(m_{kl}^{q})^{1/\theta}$.
\end{itemize}
\end{frame}
\note{\textbf{Talk track, 60 seconds.} For each adjacent grid pair, I select the path that minimizes observed morning minutes on the directed raw-road network. The figure traces one actual Beijing pair from the selected raw-road pieces and endpoint access links to observed physical time. Observed and free-flow speeds are applied to the same selected path. The model then maps physical minutes into $t_{kl}=m_{kl}^{1/\theta}$ before constructing route and bilateral costs.}

\begin{frame}[label=estsolve]{Estimation: Invert, Anchor, and Solve}
\begin{itemize}\small\setlength\itemsep{0.55em}
  \item \textbf{Invert fundamentals.} Recover $(\bar u_i^\theta,\bar A_j^\theta,\chi)$ from observed residence and workplace marginals at $m_{kl}^{obs}$. \app{appodfit}
  \item \textbf{Anchor baseline.} Hold $m_{kl}=m_{kl}^{obs}$, solve the baseline equilibrium, and obtain model-implied $\Xi_{kl}^{0}$.
  \item \textbf{Solve the tidal-lane fixed point in physical minutes:}
  \vspace{-0.45em}\begin{equation}\label{eq:anchor}
    t_{kl}^{(r)}=\big(m_{kl}^{(r)}\big)^{1/\theta},\qquad
    \widetilde m_{kl}^{(r+1)}=\max\!\left\{
    m_{kl}^{ff},\;
    m_{kl}^{obs}
    \left[
    \frac{\Xi_{kl}^{(r)}}{\Xi_{kl}^{0}}
    \frac{n_{kl}^{0}}{n_{kl}^{cf}}
    \right]^{\lambda}
    \right\}.
  \end{equation}
  \vspace{-0.65em}
  \item \textbf{Iterate.} Recompute $(t,\tau,\Xi,L^R,L^F)$, damp $m$, and repeat; ST fixes $(L^R,L^F)$, whereas LT re-solves them.
\end{itemize}
\end{frame}
\note{\textbf{Talk track, 55 seconds.} The inversion recovers fundamentals that rationalize the observed marginals. Holding observed physical minutes fixed produces the baseline model link load. A tidal-lane counterfactual changes directional lanes, then iterates among physical minutes, model link traffic, and location choices. Physical minutes respond to traffic per lane with elasticity $\lambda=1.1$; iceberg cost is obtained afterward as $t=m^{1/\theta}$. The link-load ratio maps the model's proportional traffic response into the congestion equation.}

\begin{frame}[label=fundamentals]{Baseline Fundamentals (inverted)}
\begin{columns}[T,onlytextwidth]
\column{0.5\textwidth}
\centering \scriptsize Residence amenity $\bar u_i$\\
\includegraphics[width=\linewidth,height=0.52\textheight,keepaspectratio]{fund_ubar.png}
\column{0.5\textwidth}
\centering \scriptsize Workplace productivity $\bar A_j$\\
\includegraphics[width=\linewidth,height=0.52\textheight,keepaspectratio]{fund_abar.png}
\end{columns}
\vspace{0.1em}
\begin{itemize}\footnotesize
  \item $\bar u,\bar A$ are inverted to match observed residence and workplace shares at observed travel times.
  \item $\mathrm{corr}(\bar A_j,L_j^F)=\good{0.67}$ vs.\ $\mathrm{corr}(\bar u_i,L_i^R)=0.30$: jobs pinned by productivity, residence is access-driven.
\end{itemize}
\end{frame}
\note{\textbf{Talk track, 35 seconds.} These maps are recovered objects, not directly observed amenities or productivity. The inversion matches baseline residence and workplace marginals at observed directed times. Use the maps to establish the spatial fundamentals that remain fixed in the tidal-lane counterfactuals.}

% =====================================================================
\section{Counterfactuals}
% =====================================================================
\begin{frame}{What Symmetry Misses}
\framesubtitle{For each reciprocal road pair, replace both directed travel times with their mean, then re-estimate the model on the symmetric network.}
\begin{itemize}\small\setlength\itemsep{0.25em}
  \item \textbf{Physical time.} More commuters travel in the slower direction, so reciprocal averaging understates experienced AM travel time.
  \item \textbf{Recovered fundamentals.} Symmetry yields lower central residential fundamentals by attenuating peak travel time from outlying residential areas.
\end{itemize}
\vspace{0.10em}
\begin{columns}[T,onlytextwidth]
\column{0.40\textwidth}
\centering
{\scriptsize\textbf{AM exposure and physical-time bias}\par}
\vspace{0.35em}
{\scriptsize
\hspace{0.04\linewidth}%
\begin{tabularx}{0.93\linewidth}{@{}Xr@{}}
\toprule
\textbf{Statistic} & \textbf{Share / change} \\
\midrule
Slower direction & \textbf{56.3\%} \\
Faster direction & 43.7\% \\
\addlinespace[0.25em]
Fixed passages & \good{\textbf{$-5.0\%$}} \\
Full re-estimation & \good{\textbf{$-5.5\%$}} \\
\bottomrule
\end{tabularx}
}
\vspace{0.45em}
{\tiny Changes are relative to the directed baseline.}

\column{0.58\textwidth}
\centering
{\scriptsize\textbf{Revision in recovered fundamentals}\par}
\vspace{0.18em}
\includegraphics[width=\linewidth,height=0.43\textheight,keepaspectratio]{symmetry_fundamental_relative_change.pdf}
\end{columns}
\end{frame}
\note{\textbf{Talk track, 35 seconds.} Arithmetic symmetrization leaves the unweighted mean link time unchanged. The slower direction carries 56.3 percent of baseline model-implied passages, versus 43.7 percent in the faster direction. Equalizing reciprocal times lowers passage-weighted physical time by 5.0 percent with passages fixed and by 5.5 percent after routes, OD flows, fundamentals, and equilibrium are re-estimated. On the right, the symmetric-minus-directed residential fundamental is negative near the center and rises with distance. Symmetry yields lower central residential fundamentals by attenuating peak travel time from outlying residential areas. Because residence and workplace fundamentals are jointly re-inverted, the profile is a spatial-attribution result rather than a one-equation causal decomposition.}

\begin{frame}{Where the Tidal Policy Lands}
\framesubtitle{Top-100 corridors by TidalCandidateScore (TCS).}
\begin{itemize}\footnotesize\setlength\itemsep{0.35em}
  \item \textbf{Direction.} For each reciprocal pair, label the slower AM direction $s$ and the faster direction $f$.
  \item \textbf{Reallocation.} Move one lane from $f$ to $s$, holding total road width fixed: $n_s'=n_s+1$ and $n_f'=n_f-1$.
  \item \textbf{Targeting.} Rank reciprocal pairs by slow-minus-fast congestion-delay mass:
  \vspace{-0.50em}\[
    \mathrm{TCS}=\Xi_s(m_s^{obs}-m_s^{ff})-\Xi_f(m_f^{obs}-m_f^{ff}).
  \]
\end{itemize}
\vspace{-0.90em}
\noindent
\begin{minipage}[b]{0.57\textwidth}
  \centering
  \includegraphics[page=23,trim=220 55 30 65,clip,width=0.72\linewidth,keepaspectratio]{tidal_advisor_revision_round2_20260703.pdf}
\end{minipage}\hfill
\begin{minipage}[b]{0.38\textwidth}
  \centering
  \includegraphics[width=0.50\linewidth,keepaspectratio]{map_treated_overlay.png}\par
  \vspace{-0.35em}
  {\scriptsize The top 100 pairs are treated. Red edges receive the lane.\par}
\end{minipage}
\end{frame}
\note{\textbf{Talk track, 60 seconds.} The intervention is capacity reallocation, not capacity expansion. One lane moves from the faster AM direction to the slower direction on each selected reciprocal corridor. TCS ranks corridors by the difference in baseline flow times excess traversal time. The map identifies the lane-receiving directions.}

\begin{frame}{Short Run vs.\ Long Run}
\begin{itemize}\small\setlength\itemsep{0.48em}
  \item \textbf{Short run (ST):} pin residence $L^R$ and employment $L^F$ at baseline; route traffic, edge times $t_{kl}$, and bilateral costs $\tau_{ij}$ re-equilibrate.
  \item \textbf{Long run (LT):} allow $L^R$, $L^F$, and OD shares to sort, with congestion feedback.
  \item $s_{ij}^{h}$ is the OD commuter share, $\tau_{ij}^{h}$ is bilateral commuting cost, and $\bar\tau^{h}=\sum_{ij}s_{ij}^{h}\tau_{ij}^{h}$ for $h\in\{0,ST,LT\}$.
\end{itemize}
\vspace{0.10em}
{\footnotesize
\begin{equation*}
\Delta\bar\tau^{LT}=
\underbrace{\textstyle\sum_{ij}s_{ij}^0\left(\tau_{ij}^{ST}-\tau_{ij}^0\right)}_{\text{ST direct}}
+\underbrace{\textstyle\sum_{ij}s_{ij}^0\left(\tau_{ij}^{LT}-\tau_{ij}^{ST}\right)}_{\text{GE wedge}}
+\underbrace{\textstyle\sum_{ij}\left(s_{ij}^{LT}-s_{ij}^0\right)\tau_{ij}^{LT}}_{\text{sorting}}.
\end{equation*}
}
\vspace{-0.55em}
\begin{itemize}\small
  \item \textbf{Decomposition terms:}
  \begin{itemize}\footnotesize\setlength\itemsep{0.18em}
    \item \textbf{ST direct:} baseline OD shares, baseline to short-run congestion.
    \item \textbf{GE wedge:} baseline OD shares, short-run to long-run congestion feedback.
    \item \textbf{Sorting:} OD shares move from $s^0$ to $s^{LT}$, evaluated at LT costs.
  \end{itemize}
\end{itemize}
\end{frame}
\note{\textbf{Talk track, 55 seconds.} The short run pins residence and employment at baseline while route traffic, edge times, and bilateral costs re-equilibrate. The long run allows residence, employment, and OD shares to sort with congestion feedback. The decomposition separates the baseline-to-short-run direct effect, the short-run-to-long-run congestion wedge, and sorting across OD pairs evaluated at long-run costs.}

\begin{frame}[label=result2]{Short-Run Time Gains Are Reversed by Long-Run Reallocation}
\framesubtitle{Incremental components and LT totals relative to the directed baseline.}
\centering
\vspace{0.10em}
\scriptsize
\renewcommand{\arraystretch}{1.06}
\setlength{\tabcolsep}{14pt}
\begin{tabular}{p{0.62\textwidth}r}
\toprule
Channel & Change \\
\midrule
\multicolumn{2}{l}{\textbf{A. Total physical commuting time (hours per AM peak)}} \\
\hspace{1.2em}ST equilibrium (fixed sorting) & \good{$-1{,}031$} \\
\hspace{1.2em}$+$ LT link-time feedback & \bad{$+650$} \\
\hspace{1.2em}$+$ LT link-passage reallocation & \bad{$+733$} \\
\cmidrule(lr){1-2}
\hspace{1.2em}\textbf{LT total (relative to baseline)} & \textbf{\bad{$+352$}} \\
\addlinespace[0.45em]
\multicolumn{2}{l}{\textbf{B. Expected-utility welfare (percentage points)}} \\
\hspace{1.2em}ST direct & \good{$+0.8714$} \\
\hspace{1.2em}$+$ GE wedge & \bad{$-0.5470$} \\
\hspace{1.2em}$+$ Sorting & \bad{$-0.0001$} \\
\cmidrule(lr){1-2}
\hspace{1.2em}\textbf{LT total (relative to baseline)} & \textbf{\good{$+0.3243$}} \\
\bottomrule
\end{tabular}

\vspace{0.34em}
{\tiny\color{gray!80}\begin{minipage}{0.90\textwidth}
Physical time is computed from additive link minutes and equilibrium link passages; negative changes are time savings. The welfare rows follow Slide 22: ST direct, GE wedge, and sorting. Sorting combines workplace productivity and residential crowding. \app{appphysicaltime}
\end{minipage}\par}

\vspace{0.42em}
\begin{minipage}{0.92\textwidth}
\centering\footnotesize
Long-run feedback reverses the physical-time gain, while welfare remains above baseline.
\end{minipage}
\end{frame}
\note{\textbf{Talk track, 110 seconds.} Panel A decomposes physical time. The short-run equilibrium saves 1,031 hours in one morning peak. Long-run link-time feedback offsets 650 hours, and link-passage reallocation offsets another 733 hours. The explicit LT total is 352 additional hours relative to baseline. Panel B follows the three-term decomposition on the preceding slide. The ST direct term raises expected utility by 0.8714 percent. The GE wedge subtracts 0.5470 percentage point. Sorting combines the workplace-productivity and residential-crowding channels and subtracts 0.0001 percentage point. These three terms add exactly to the long-run welfare gain of 0.3243 percent.}

\begin{frame}{Long-Run Sorting}
\begin{columns}[T,onlytextwidth]
\column{0.35\textwidth}
\begin{itemize}\footnotesize
  \item Totals are fixed at $\bar L$; maps report demeaned share changes, not citywide growth.
  \item Positive cells gain share relative to baseline; black circles mark treated-corridor endpoint nodes.
  \item \textbf{Residential shares} rise near treated endpoints because lower OD costs raise residential market access.
  \item \textbf{Employment gains} are concentrated in a smaller set of red cells.
\end{itemize}
\column{0.63\textwidth}
\centering
{\scriptsize
\begin{tabular}{lcccc}
\toprule
 & \multicolumn{2}{c}{Treated endpoints} & \multicolumn{2}{c}{Rest} \\
 & $\ell^R$ & $\ell^F$ & $\ell^R$ & $\ell^F$ \\
\midrule
Mean $\Delta\log$ & \good{$+0.017$} & $-0.006$ & \bad{$-0.012$} & \bad{$-0.016$} \\
\bottomrule
\end{tabular}
}

\vspace{1em}
\begin{minipage}[t]{0.49\linewidth}
\centering \scriptsize demeaned $\Delta\log\ell^R$\\
\includegraphics[width=\linewidth,height=0.45\textheight,keepaspectratio]{sort_lR.png}
\end{minipage}\hfill
\begin{minipage}[t]{0.49\linewidth}
\centering \scriptsize demeaned $\Delta\log\ell^F$\\
\includegraphics[width=\linewidth,height=0.45\textheight,keepaspectratio]{sort_lF.png}
\end{minipage}
\end{columns}
\end{frame}

\begin{frame}[label=result4]{Transportation Policy Comparison}
\framesubtitle{Endogenous GE under anchored congestion and long-run sorting.}
\begin{center}\scriptsize
\begin{tabular}{lcccc}
\toprule
Policy & ST $\bar W$ & LT $\bar W$ & net lane-km & cost/lane-km \\
\midrule
Tidal reallocation & $+0.87\%$ & $+0.32\%$ & \good{\textbf{0}} & near zero capex \\
Road Expansion & $+1.26\%$ & $+0.44\%$ & 1{,}140 & $\approx$ CNY 5m \\
\bottomrule
\end{tabular}
\end{center}
\vspace{0.3em}
\begin{itemize}\footnotesize
  \item Tidal reallocation lifts LT welfare $+0.32\%$ at zero net lane-km.
  \item Delivering comparable long-run welfare by building capacity takes about $1{,}140$ lane-km: Road Expansion adds lanes where baseline flow$\times$delay is highest and reaches $+0.44\%$.
  \item Urban arterial widening runs on the order of CNY 5m per lane-km (roughly CNY 30m/km for a six-lane arterial), so the expansion benchmark is order CNY 6bn; tidal reallocation avoids this outlay. \app{apppolicy}
\end{itemize}
\end{frame}
\note{\textbf{Talk track, 45 seconds.} The comparison is limited to welfare, capacity, and construction cost because the two policies do not share the same treated-link support. Tidal reallocation changes the direction in which existing capacity is deployed; the expansion benchmark builds about 1,140 lane-kilometers. The preceding slide reports citywide physical time separately from expected utility.}

\begin{frame}{Conclusion}
\begin{enumerate}
  \item \textbf{Directional congestion is a widespread urban feature.}
  \begin{itemize}\footnotesize
    \item Across the top 100 congested cities, flow-weighted median \SAR\ is \good{0.878}; the same asymmetry persists on Beijing road segments.
  \end{itemize}
  \item \textbf{Symmetry misses commuter exposure and spatial attribution.}
  \begin{itemize}\footnotesize
    \item \good{$56.3\%$} of AM passages take the slower direction; symmetry understates travel time by \good{about $5\%$}.
    \item Symmetry understates peripheral peak travel time and, in turn, central residential fundamentals.
  \end{itemize}
  \item \textbf{Tidal lanes deliver welfare gains through directed capacity reallocation.}
  \begin{itemize}\footnotesize
    \item ST total time saved: \good{$1{,}031$ hours per AM peak}; after LT reallocation, total time rises by \bad{$352$ hours} and $\bar W$ remains \good{$+0.32\%$}.
    \item Supply-side policies with comparable LT welfare require about \bad{$1{,}100$ lane-km}; tidal reallocation uses \good{zero net lane-km}.
  \end{itemize}
\end{enumerate}
\end{frame}
\note{\textbf{Talk track, 25 seconds.} Directional congestion is widespread. The more congested direction carries 56.3 percent of AM link passages, so symmetry understates experienced travel time by about 5 percent. It also understates peripheral peak travel time and, in turn, central residential fundamentals. Tidal reallocation delivers short-run time savings and positive long-run welfare with zero net lane-kilometers.}

% =====================================================================
\appendix
\section{Appendix}
% =====================================================================
\begin{frame}[label=appmeasurementrouting]{Appendix: Selecting the Shortest Directed Road Path}
\begin{columns}[T,onlytextwidth]
\column{0.54\textwidth}
\begin{enumerate}\small\setlength\itemsep{0.75em}
  \item \textbf{Build endpoint candidates.} Connect each grid centroid to at most three nearby directed road edges within 2.5 km.
  \item \textbf{Solve every candidate pairing.} For each origin-edge and destination-edge pair, Dijkstra finds the directed road path with the lowest observed AM minutes. Access minutes are added at both endpoints.
  \item \textbf{Select the global minimum.} Retain the candidate path with the lowest total minutes:
\end{enumerate}
\vspace{-0.25em}
\begin{equation*}
\begin{aligned}
  \text{total minutes}={}&\text{origin access}+\text{road-path minutes}\\
  &+\text{destination access}.
\end{aligned}
\end{equation*}
\vspace{-0.35em}
{\footnotesize\textbf{Direction matters.} The $k\to l$ and $l\to k$ paths are solved separately.\par}

\column{0.43\textwidth}
\centering
\small\textbf{Selected path for one adjacent grid pair}\par
\vspace{0.35em}
\includegraphics[width=\linewidth,height=0.51\textheight,keepaspectratio]{figs/route_examples/route_example_01.png}\par
\vspace{-0.25em}
{\scriptsize Blue: selected directed road path.\par
Orange: endpoint access connectors.\par}
\end{columns}
\navback{linkmeasurement}
\end{frame}

\begin{frame}[label=approute]{Appendix: Route Costs and Bilateral Commuting Costs}
\begin{itemize}\small
  \item Define the directed adjacency matrix $A=[a_{kl}]$ by
  \vspace{-0.5em}\begin{equation}
    a_{kl} = t_{kl}^{-\theta}
  \end{equation}
  \item Summing over all feasible route lengths gives
  \vspace{-0.5em}\begin{equation}
    B = \sum_{K=0}^{\infty}A^K = (I-A)^{-1}
  \end{equation}
  \item The bilateral commuting-cost aggregator can be written as
  \vspace{-0.5em}\begin{equation}\label{eq:inverse}
    \tau_{ij} = B_{ij}^{-1/\theta}
  \end{equation}
  \vspace{-0.5em}
  \item The current mainline uses common route rescaling $\zeta=1$, so $A_{kl}=1/m_{kl}$ exactly. The solver retains an optional common rescaling for cases in which the spectral-radius condition would otherwise fail.
\end{itemize}
\navback{bilateral}
\end{frame}

\begin{frame}[label=appphysicaltime]{Appendix: Expected Physical Travel Time}
\scriptsize
\begin{itemize}\setlength\itemsep{0.45em}
  \item Physical minutes add along route $r$:
  \vspace{-0.35em}
  \[
    M_{ij,r}=\sum_{(k,l)\in r}m_{kl}.
  \]
  \vspace{-0.75em}
  \item Iceberg link costs compound, $C_{ij,r}=\prod_{(k,l)\in r}t_{kl}$, and the conditional route probability is
  \vspace{-0.35em}
  \[
    \rho_{ij,r}=\frac{C_{ij,r}^{-\theta}}{\sum_{s\in\mathcal R_{ij}}C_{ij,s}^{-\theta}},
    \qquad t_{kl}=m_{kl}^{1/\theta}.
  \]
  \vspace{-0.75em}
  \item Expected OD minutes and the citywide mean are
  \vspace{-0.35em}
  \[
    M_{ij}=\sum_r\rho_{ij,r}M_{ij,r},
    \qquad
    \bar M=\frac{1}{\bar L}\sum_{i,j}L_{ij}M_{ij}
    =\frac{1}{\bar L}\sum_{k,l}\Xi_{kl}m_{kl}.
  \]
  \vspace{-0.75em}
  \item No transformation of $\tau_{ij}$ produces additive route minutes. The implementation computes $\bar M$ from equilibrium link passages and $m_{kl}$.
\end{itemize}
\navback{result2}
\end{frame}

\begin{frame}[label=apptreatedtime]{Appendix: Treated-Link Time Incidence}
\framesubtitle{A local engineering incidence measure, distinct from citywide experienced travel time.}
\footnotesize
\textbf{Baseline-passage accounting}\par
\vspace{0.10em}
\begin{align*}
\text{Gross savings}^{h}
  &=\sum_{e\in\mathcal E^*}\max\{m_e^0-m_e^h,0\}\,\Xi_e^0,\\
\text{Offsetting losses}^{h}
  &=-\sum_{e\in\mathcal E^*}\max\{m_e^h-m_e^0,0\}\,\Xi_e^0.
\end{align*}
\vspace{-0.45em}
\begin{center}
\renewcommand{\arraystretch}{1.18}
\begin{tabular}{lcc}
\toprule
Person-hours & ST & LT \\
\midrule
Gross savings & $1{,}385$ & $804$ \\
Offsetting losses & $-109$ & $-90$ \\
Net savings & $1{,}276$ & $714$ \\
\bottomrule
\end{tabular}
\end{center}
\vspace{0.15em}
\begin{itemize}\scriptsize\setlength\itemsep{0.25em}
  \item $\mathcal E^*$ contains the 200 treated directed links; $\Xi_e^0$ fixes baseline passage weights.
\end{itemize}
\navback{result2}
\end{frame}

\begin{frame}[label=appfund]{Appendix: Estimation Algorithm}
\small\textbf{1. Baseline inversion}\par
\vspace{0.08em}
\begin{itemize}\footnotesize\setlength\itemsep{0.22em}
  \item Set $t_{kl}^{obs}=(m_{kl}^{obs})^{1/\theta}$ and recover $(\bar u_i^\theta,\bar A_j^\theta,\chi)$ by alternating updates that match $(L_i^R,L_j^F)$.
  \item Holding $m=m^{obs}$, solve route and location choices to construct $\Xi_{kl}^{0}$; gravity determines model-implied $L_{ij}$.
\end{itemize}
\vspace{0.10em}
\small\textbf{2. Counterfactual fixed point}\par
\vspace{0.04em}
\begin{itemize}\footnotesize\setlength\itemsep{0.18em}
  \item Replace $n^0$ with $n^{cf}$, initialize $m^{(0)}=m^{obs}$, and hold fundamentals and calibrated parameters fixed.
\end{itemize}
\vspace{-0.62em}
{\footnotesize\begin{align*}
  t_{kl}^{(r)}&=\big(m_{kl}^{(r)}\big)^{1/\theta},
  &\Xi_{kl}^{(r)}&=\Xi_{kl}\!\left(t^{(r)},L^{R,(r)},L^{F,(r)}\right),\\[-0.15em]
  \widetilde m_{kl}^{(r+1)}
  &=\max\!\left\{m_{kl}^{ff},\;
  m_{kl}^{obs}\left[
  \frac{\Xi_{kl}^{(r)}}{\Xi_{kl}^{0}}
  \frac{n_{kl}^{0}}{n_{kl}^{cf}}
  \right]^{\lambda}\right\},
  &\log m_{kl}^{(r+1)}
  &=(1-\omega)\log m_{kl}^{(r)}+\omega\log\widetilde m_{kl}^{(r+1)}.
\end{align*}}
\vspace{-0.65em}
\begin{itemize}\footnotesize\setlength\itemsep{0.18em}
  \item Iterate until $m$ converges. ST pins $(L^R,L^F)$; LT re-solves them. Mainline values are $\lambda=1.1$ and $\omega=0.25$.
\end{itemize}
\vspace{-0.03em}
{\scriptsize\color{gray!80} Implied traffic-index interpretation: $q_{kl}^{0}=n_{kl}^{0}(m_{kl}^{obs}/m_{kl}^{ff})^{1/\lambda}$ and $\widehat q_{kl}^{(r)}=\widehat\Xi_{kl}^{(r)}$; the code eliminates $q^0$ algebraically.}
\navback{fundamentals}
\end{frame}

\begin{frame}[label=appodfit]{Appendix: Bilateral OD Fit beyond Marginals}
\begin{minipage}{0.94\textwidth}
\footnotesize
\textbf{Calculation.} Observed and predicted log OD flows are residualized separately on origin and destination fixed effects; $\rho$ is their pairwise residual correlation.
\end{minipage}

\vspace{0.20em}
\centering
\includegraphics[
  width=0.78\textwidth,
  height=0.58\textheight,
  keepaspectratio
]{../../data_work/outputs/od_validation_aa_style_gate5_bpr_pos05_v1/figures/figure_od_validation_structural_central99.pdf}

\vspace{-0.05em}
{\tiny\color{gray!75!black} Axes omit the outer 0.5\% of each residual distribution; the correlation uses the full pairwise sample.\par}

\vspace{0.22em}
\begin{minipage}{0.94\textwidth}
\footnotesize
\textbf{Interpretation.} With residence and workplace scale absorbed, $\rho=\good{\mathbf{0.413}}$ evaluates bilateral variation implied by directed network costs and gravity; it is not a unit-slope calibration of flow levels.
\end{minipage}
\navback{estsolve}
\end{frame}
\note{\textbf{Q\&A only.} The fundamentals inversion matches observed residence and workplace marginals, not the bilateral OD matrix. Observed and predicted log flows are residualized separately on origin and destination fixed effects. Their residual correlation is 0.413 across 83,030 positive observed pairs. The statistic captures bilateral commuting variation beyond origin and destination scale; it does not imply one-for-one calibration of residual flow magnitudes.}

\begin{frame}[label=appduncan]{Appendix: Jobs--Housing Dissimilarity Benchmarks}
\begin{itemize}\footnotesize
  \item \textbf{Cross-city.} Beijing's township-level jobs--housing dissimilarity is $\good{0.20}$, the highest among the top 100 cities (median $0.09$).
\end{itemize}
\vspace{0.15em}
\centering
\includegraphics[width=0.96\textwidth,height=0.59\textheight,keepaspectratio]{app_duncan_cross_city.png}
\navback{jobhousing}
\end{frame}

\begin{frame}[label=appparams]{Appendix: Parameter Interpretation}
\centering
\vspace{0.25em}
\scriptsize
\renewcommand{\arraystretch}{1.32}
\begin{tabularx}{0.94\textwidth}{@{}p{0.10\textwidth}p{0.25\textwidth}X@{}}
\toprule
Parameter & Economic role & Direct margin in the model \\
\midrule
$\theta$ & Fr\'echet dispersion & Governs sensitivity across residence, workplace, and route alternatives; also maps $m_{kl}$ into $t_{kl}=m_{kl}^{1/\theta}$. \\
$\alpha>0$ & Workplace agglomeration & Raises productivity with employment density: $w_j=\bar A_j(L_j^F)^\alpha$. \\
$\beta<0$ & Residential crowding & Lowers residential amenity with resident density: $u_i=\bar u_i(L_i^R)^\beta$. \\
$\lambda$ & Congestion sensitivity & A higher $\lambda$ makes travel time more responsive to traffic per lane. \\
\bottomrule
\end{tabularx}
\vspace{0.55em}

\begin{minipage}{0.92\textwidth}
\footnotesize Sorting combines accessibility, workplace agglomeration, and residential crowding. The current specification calibrates all four parameters externally.
\end{minipage}
\navback{estparams}
\end{frame}

\begin{frame}[label=appsens]{Appendix: Sensitivity Summary}
\framesubtitle{Final mainline rerun, TCS top-100, primitive $\bar W$, $\lambda=1.1$.}
\vspace{0.22em}
\centering
\footnotesize
\renewcommand{\arraystretch}{1.16}
\begin{tabular}{@{}cc@{\hspace{2.6em}}cc@{}}
\toprule
\multicolumn{2}{c}{\textbf{Vary $\alpha$; hold $\beta=-0.10$}} &
\multicolumn{2}{c}{\textbf{Vary $\beta$; hold $\alpha=+0.05$}} \\
\midrule
$\alpha$ & LT $\bar W$ & $\beta$ & LT $\bar W$ \\
\midrule
$-0.12$ & $+0.37\%$ & $\mathbf{-0.10}$ & $\mathbf{+0.32\%}$ \\
$0.00$ & $+0.34\%$ & $0.00$ & $+0.31\%$ \\
$+0.03$ & $+0.33\%$ & $+0.05$ & $+0.28\%$ \\
$\mathbf{+0.05}$ & $\mathbf{+0.32\%}$ & $+0.10$ & \emph{not converged} \\
$+0.10$ & $+0.29\%$ & & \\
\bottomrule
\end{tabular}
\par
\vspace{0.38em}
\raggedright
\small\textbf{Interpretation}\par
\vspace{0.08em}
\begin{itemize}\footnotesize\setlength\itemsep{0.12em}
  \item LT welfare remains positive in every converged $\alpha$--$\beta$ cell.
  \item Every cell on this page uses the current anchored specification, TCS top-100 treatment, and $\lambda=1.1$.
  \item The $\theta$ grid re-inverts fundamentals and re-solves each equilibrium. \app{appsens_theta}
  \item The asymmetry-ratio selection rule is weaker than TCS selection and can reverse the LT welfare sign.
\end{itemize}
\navback{result2}
\end{frame}

\begin{frame}[label=appsens_theta]{Appendix: Theta Sensitivity}
\framesubtitle{TCS top-100; each $\theta$ re-inverts fundamentals and re-solves baseline + tidal ST/LT.}
\scriptsize
\begin{center}
\renewcommand{\arraystretch}{1.3}
\setlength{\tabcolsep}{12pt}
\begin{tabular}{ccccc}
\toprule
$\theta$ & ST $\bar W$ & LT $\bar W$ & res.\ dissim & job dissim \\
\midrule
$3$ & $+1.99\%$ & $+0.75\%$ & $1.11$ & $1.55$ \\
$5$ & $+1.19\%$ & $+0.45\%$ & $0.97$ & $1.67$ \\
$\mathbf{6.83}$ & $\mathbf{+0.87\%}$ & $\mathbf{+0.32\%}$ & $\mathbf{0.86}$ & $\mathbf{1.79}$ \\
$8$ & $+0.74\%$ & $+0.27\%$ & $0.80$ & $1.89$ \\
$10$ & $+0.59\%$ & $+0.21\%$ & $0.70$ & $2.07$ \\
\bottomrule
\end{tabular}
\end{center}
\vspace{0.4em}
\begin{itemize}\footnotesize
  \item LT $\bar W$ stays positive across the grid; magnitudes fall as $\theta$ rises (stronger cost sensitivity).
  \item Resident dissim falls and job dissim rises with $\theta$; baseline $\bar W$ levels are not comparable across $\theta$.
\end{itemize}
\note{Mainline $\theta=6.83$ matches the tidal diagnostics row; baseline $\bar W$ levels are not comparable across $\theta$.}
\navback{appsens}
\end{frame}

\begin{frame}[label=apppolicy]{Appendix: Policy Definitions and Selected Effects}
\scriptsize
\begin{center}
\renewcommand{\arraystretch}{1.20}
\begin{tabularx}{0.96\textwidth}{@{}p{0.17\textwidth}Xccc@{}}
\toprule
Policy & Capacity rule & ST $\bar W$ & LT $\bar W$ & lane-km \\
\midrule
Tidal reallocation & TCS top-100; add one lane in the slow direction and remove one from the reciprocal fast direction. & $+0.87\%$ & $+0.32\%$ & 0 \\
Targeted expansion & Top 100 by baseline flow $\times$ excess travel time; add one lane in each direction. & $+1.26\%$ & $+0.44\%$ & 1{,}140 \\
Uniform expansion & Scale all links by $+2.4\%$ so added lane-km matches the targeted-expansion budget. & $+0.21\%$ & $+0.10\%$ & 1{,}092 \\
\bottomrule
\end{tabularx}
\end{center}
\vspace{0.35em}
\begin{itemize}\footnotesize\setlength\itemsep{0.40em}
  \item ST pins residence and workplace shares at baseline; LT re-solves residence--workplace sorting.
  \item Targeted expansion assigns capacity using baseline flow and delay; uniform expansion spreads a comparable lane-km budget over the network.
  \item The cost benchmark is roughly CNY 5m per widened urban-arterial lane-km, placing targeted expansion near CNY 6bn. Tidal reallocation has near-zero construction capex.
\end{itemize}
\navback{result4}
\end{frame}

\begin{frame}[label=resultelast]{Appendix: Targeting Tidal Lanes}
\framesubtitle{Per-link welfare elasticity from $d\bar W^\theta=dA_{ab}\phi_a\psi_b$.}
\begin{columns}[T,onlytextwidth]
\column{0.44\textwidth}
\centering
\includegraphics[width=\linewidth,height=0.66\textheight,keepaspectratio]{map_link_welfare_elast.png}
\column{0.53\textwidth}
\scriptsize
\textbf{Selection rule} (top-100):
\begin{tabular}{lc}
\toprule
rule & LT $\bar W$ \\
\midrule
\good{welfare elast.} & \good{$+0.35\%$} \\
TCS score & $+0.32\%$ \\
speed only & $+0.13\%$ \\
\bottomrule
\end{tabular}
\vspace{0.5em}
\begin{itemize}\footnotesize
  \item The local score decomposes into edge gain $\times$ upstream residence access $\times$ downstream workplace opportunity. \app{appelast}
  \item Flow matters: the speed-only rank correlation with welfare elasticity is \bad{$0.08$}; the flow$\times$delay correlation is $0.65$.
  \item Welfare-elasticity and TCS targeting select the high-flow spine; speed-only also selects slow, low-flow links.
\end{itemize}
\end{columns}
\end{frame}

\begin{frame}[label=appelast]{Appendix: Welfare-Elasticity Score}
\scriptsize
\begin{itemize}
  \item Define the directed-link accessibility matrix by $A_{ab}=t_{ab}^{-\theta}$ and the network multiplier by $B=(I-A)^{-1}$.
  \item A local gain on one directed edge uses the matrix-inverse differential:
  \[
    dB=B(dA)B.
  \]
  \item If only edge $a\to b$ changes, $dA$ has one nonzero entry, so the welfare-index numerator has the local derivative
  \[
    d\bar W^\theta=dA_{ab}\phi_a\psi_b.
  \]
  \item $dA_{ab}$ is the direct gain from one more lane; $\phi_a$ is upstream residence-access weight reaching node $a$; $\psi_b$ is downstream workplace-opportunity weight from node $b$.
  \item The score is a local ranking device. The reported LT $\bar W$ values still come from re-solving the endogenous GE counterfactual.
\end{itemize}
\navback{resultelast}
\end{frame}

\begin{frame}[label=nextsteps]{Appendix: Validation and Extensions}
\vspace{-0.30em}
{\small\textbf{Case:} Shenzhen tidal-lane application on Caitian Road. \app{appvideofacts}\par}
\vspace{0.28em}
{\small\textbf{Next step: data collection}\par}
\vspace{0.02em}
\begin{itemize}\footnotesize\setlength\itemsep{0.12em}
  \item \textbf{Coverage.} Record at least 60 minutes before and after each scheduled cutoff $c_d$ on repeated opening dates.
  \item \textbf{Outcomes.} Recover directional flow $Q_{dt}$, matched travel time $T_{dt}$, and actual opening $D_{dt}=Open_{dt}$.
\end{itemize}
\vspace{0.16em}
\noindent\makebox[\textwidth][c]{%
\begin{minipage}[t]{0.175\textwidth}
\centering
\includegraphics[width=\linewidth]{fieldvisit_crop_1.jpg}\par\vspace{0.18em}
\includegraphics[width=\linewidth]{fieldvisit_crop_2.jpg}
\end{minipage}\hspace{0.030\textwidth}%
\begin{minipage}[t]{0.175\textwidth}
\centering
\includegraphics[width=\linewidth]{fieldvisit_crop_5.jpg}\par\vspace{0.18em}
\includegraphics[width=\linewidth]{fieldvisit_crop_4.jpg}
\end{minipage}\hspace{0.030\textwidth}%
\begin{minipage}[t]{0.175\textwidth}
\centering
\includegraphics[width=\linewidth]{fieldvisit_crop_6.jpg}\par\vspace{0.18em}
\includegraphics[width=\linewidth]{fieldvisit_crop_7.jpg}
\end{minipage}%
}
\end{frame}

\begin{frame}[label=appvideofacts]{Appendix: Video Evidence Before and After Opening}
\vspace{-0.38em}
\noindent\makebox[\textwidth][c]{%
\begin{minipage}[t]{0.34\textwidth}
  \centering
  \includegraphics[width=0.74\linewidth,trim=140 50 140 90,clip]{caitian_0709_pre_tracking.jpg}\\[-0.1em]
  {\small\textbf{Pre}}
\end{minipage}\hspace{0.035\textwidth}%
\begin{minipage}[t]{0.34\textwidth}
  \centering
  \includegraphics[width=0.74\linewidth,trim=140 50 140 90,clip]{caitian_0709_late_post_tracking.jpg}\\[-0.1em]
  {\small\textbf{Post}}
\end{minipage}%
}
\vspace{0.20em}

\centering
{\scriptsize
\setlength{\tabcolsep}{0.42em}
\renewcommand{\arraystretch}{1.08}
\resizebox{0.82\textwidth}{!}{%
\begin{tabular}{@{}lrrr@{\hspace{0.60em}}rrr@{}}
\toprule
& \multicolumn{3}{c}{\textbf{Flow (vehicles/hour)}}
& \multicolumn{3}{c}{\textbf{Travel time (seconds)}} \\
\cmidrule(lr){2-4}\cmidrule(lr){5-7}
\textbf{Direction} & \textbf{Pre} & \textbf{Post} & \textbf{Change}
& \textbf{Pre} & \textbf{Post} & \textbf{Change} \\
\midrule
\good{$S\to N$ \textbf{(+1 lane)}} & 2,902 & 3,284 & \good{$+13.2\%$}
& 18.73 & 8.87 & \good{$-52.6\%$} \\
$N\to S$ & 1,484 & 1,712 & \good{$+15.4\%$}
& 3.49 & 3.90 & \bad{$+11.8\%$} \\
\bottomrule
\end{tabular}}}

\vspace{0.55em}
\raggedright
{\footnotesize\textbf{Descriptive pattern.} Flow rises in both directions. Median travel time falls only for $S\to N$ (+1 lane).\par}

\vfill
{\fontsize{5.0}{5.25}\selectfont\color{gray!75!black}\raggedright
\textbf{Data notes.} Pre/Post windows are 0:00--2:45 and 33:30--48:30; deployment is excluded. Flow is normalized to vehicles/hour, and travel time is the median matched gate time.\par}

\par\vspace{0.18em}\noindent\hfill
\hyperlink{motivation}{\beamerreturnbutton{motivation}}\;
\hyperlink{nextsteps}{\beamerreturnbutton{validation}}\;
\hyperlink{appvideordit}{\beamergotobutton{design}}
\end{frame}
\note{\textbf{Q\&A only.} The 0709 video comparison documents directed flow and matched gate time before and after opening. The evidence is descriptive because it comes from one day and two separated windows. The repeated-date RDiT design on the next slide defines the causal estimand.}

\begin{frame}[label=appvideordit]{Appendix: Opening-Time RDiT Design}
\vspace{-0.48em}
\small
\textbf{Design.}
At the scheduled cutoff $c_d$, $S\to N$ capacity shifts from three to four lanes. Actual operation $D_{dt}=Open_{dt}$ is coded from video; $Y_{dt}\in\{Q_{dt},T_{dt}\}$ measures directional flow or matched travel time.

\vspace{0.10em}
{\centering
$\displaystyle r_{dt}=t-c_d,\qquad
Z_{dt}=\mathbf{1}\{r_{dt}\geq0\},\qquad
D_{dt}=Open_{dt}. $\par}

\vspace{0.16em}
{\centering
$\displaystyle
Y_{dt}=\alpha_d+\tau_YD_{dt}
+\beta_-r_{dt}+\beta_+D_{dt}r_{dt}+\varepsilon_{dt},
\qquad |r_{dt}|\leq h.$\par}

\vspace{0.16em}
\textbf{Identification}\par
\begin{itemize}\footnotesize\setlength\itemsep{0em}
  \item \textbf{Predetermined cutoff.} The scheduled opening time is set before minute-level traffic shocks.
  \item \textbf{Continuity.} Absent opening, directional outcomes evolve smoothly through $c_d$.
\end{itemize}
\vspace{-0.05em}
\textbf{Estimation}\par
\begin{itemize}\footnotesize\setlength\itemsep{0em}
  \item Repeated dates enter as five-minute bins with date fixed effects and date-clustered inference.
  \item $h$ is the symmetric time bandwidth around $c_d$: only observations within $h$ minutes on either side enter. Deployment is excluded; robustness checks vary the bin width and $h$ and test placebo cutoffs.
\end{itemize}

\vspace{0em}
{\footnotesize\textbf{Estimand.} $\tau_Y$ is the local effect of actual opening. If operation is delayed, $Z_{dt}$ instruments $D_{dt}$ in a fuzzy RDiT.}

\vfill
{\tiny\color{gray!75!black} Current 0709 estimates are descriptive. Formal RDiT estimation requires repeated opening dates with continuous coverage around the cutoff.}

\par\vspace{0.18em}\noindent\hfill
\hyperlink{appvideofacts}{\beamerreturnbutton{evidence}}\;
\hyperlink{nextsteps}{\beamerreturnbutton{validation}}
\end{frame}
\note{\textbf{Q\&A only.} Scheduled opening time supplies the cutoff, while actual lane operation is coded from the video. The design requires repeated dates, continuity through the cutoff absent opening, and date-clustered inference. If actual operation differs from the schedule, the cutoff indicator instruments actual opening in a fuzzy RDiT.}

\nobibliography{references}

\end{document}
